\subsection{\texorpdfstring{The Grothendieck Group \(R_{K}(A)\)}{The Grothendieck Group R\_\{K\}(A)}}

Let \(\mathrm{K}\) be a commutative field and \(\mathrm{A}\) a \(\mathrm{K}\) -algebra. The set of classes of A-modules that are finite-dimensional vector spaces over \(\mathrm{K}\) is hereditary. The corresponding Grothendieck group is denoted by \(\mathrm{R}_{\mathrm{K}}(\mathrm{A})\) . It is a free \(\mathbf{Z}\) -module with basis the family \(([\mathrm{S}])_{\mathrm{S}\in \mathcal{S}}\) , where \(\mathcal{S}\) is the set of classes of simple A-modules that are finite-dimensional over \(\mathrm{K}\) . There exists a homomorphism

\[
\dim : \mathrm{R} _ {\mathrm{K}} (\mathrm{A}) \longrightarrow \mathbf {Z}
\]

characterized by \(\dim([E]) = [E : K]\) for every A-module E that is finite-dimensional over K. When A = K, it is an isomorphism. The submonoid of effective elements is denoted by \(\mathrm{R}_{\mathrm{K}}(\mathrm{A})^{+}\) .

Lemma 2. --- Let M and M' be A-modules that are finite-dimensional vector spaces over K. The supports (VIII, p.~190) of M and M' are disjoint if and only if there exists an \(a \in A\) such that \(a_{M} = 0_{M}\) and \(a_{M'} = 1_{M'}\) .

Suppose that there exists an \(a \in \mathrm{A}\) such that \(a_{\mathrm{M}} = 0_{\mathrm{M}}\) and \(a_{\mathrm{M}'} = 1_{\mathrm{M}'}\) . Let S be a simple A-module. If \(\operatorname{cl}(S)\) belongs to the support \(\mathcal{S}_{\mathrm{M}}\) of M, then the A-module S is isomorphic to one of the quotients of a Jordan-Hölder series of M and we have \(a_{\mathrm{S}} = 0_{\mathrm{S}}\) . Likewise, if \(\operatorname{cl}(S)\) belongs to the support \(\mathcal{S}_{\mathrm{M}'}\) of \(\mathrm{M}'\) , then we have \(a_{\mathrm{S}} = 1_{\mathrm{S}}\) . It follows that \(\mathcal{S}_{\mathrm{M}}\) and \(\mathcal{S}_{\mathrm{M}'}\) are disjoint.

Conversely, suppose that the sets \(S_{M}\) and \(S_{M'}\) are disjoint. They are finite because M and \(M'\) are finite-dimensional over K. Every simple A-module S whose class belongs to \(S_{M} \cup S_{M'}\) is finite-dimensional over K and a fortiori over the field \(\operatorname{End}_{\mathrm{A}}(\mathrm{S})\) . By Corollary 1 of Proposition 4 (VIII, p.~83), there exists an element \(b \in A\) such that we have \(b_{S} = 0_{S}\) (resp. \(b_{S} = 1_{S}\) ) for every simple A-module S whose class belongs to \(S_{M}\) (resp. \(S_{M'}\) ). Let \((\mathrm{M}_{i})_{0 \leqslant i \leqslant n}\) be a Jordan--Hölder series of M. By the above, we have \(bM_{i} \subset M_{i+1}\) for \(0 \leqslant i < n\) and therefore \((b^{n})_{\mathrm{M}} = 0_{\mathrm{M}}\) . The existence of a natural number m such that \(((1-b)^{m})_{\mathrm{M}^{\prime}}=0_{\mathrm{M}^{\prime}}\) is proved in the same way. Set \(\mathrm{P}(\mathrm{X})=1-(1-\mathrm{X}^{n})^{m}\) and \(a=\mathrm{P}(b)\) . The polynomial \(\mathrm{P}(\mathrm{X})\) is a multiple of \(X^{n}\) , so we have \(a_{M}=0_{M}\) , while the polynomial \(1-\mathrm{P}(\mathrm{X})\) is a multiple of \((1-\mathrm{X})^{m}\) , so we have \(a_{M^{\prime}}=1_{M^{\prime}}\) . This concludes the proof.

Let \(\mathrm{L}\) be an extension of \(\mathrm{K}\) . If \(\mathrm{M}\) is an A-module that is finite-dimensional over \(\mathrm{K}\) , then \(\mathrm{M}_{(\mathrm{L})}\) is an \(\mathrm{A}_{(\mathrm{L})}\) -module that is finite-dimensional over \(\mathrm{L}\) . Moreover, for every exact sequence

\[
0 \longrightarrow \mathrm{M} ^ {\prime} \longrightarrow \mathrm{M} \longrightarrow \mathrm{M} ^ {\prime \prime} \longrightarrow 0
\]

of A-modules, the sequence of \(A_{(L)}\) -modules

\[
0 \longrightarrow \mathrm{M} _ {(\mathrm{L})} ^ {\prime} \longrightarrow \mathrm{M} _ {(\mathrm{L})} \longrightarrow \mathrm{M} _ {(\mathrm{L})} ^ {\prime \prime} \longrightarrow 0
\]

deduced from it by extension of scalars is exact (II, §7, No.~7, p.~308, Proposition 14). Hence there exists a unique ring homomorphism \(u: \mathrm{R}_{\mathrm{K}}(\mathrm{A}) \to \mathrm{R}_{\mathrm{L}}(\mathrm{A}_{(\mathrm{L})})\) such that \(u([\mathrm{M}]) = [\mathrm{M}_{(\mathrm{L})}]\) for every A-module M that is finite-dimensional over K (No.~5).

THEOREM 1. --- The homomorphism \(u: \mathrm{R}_{\mathrm{K}}(\mathrm{A}) \to \mathrm{R}_{\mathrm{L}}(\mathrm{A}_{(\mathrm{L})})\) defined above is injective. Let \(\xi\) be an element of \(\mathrm{R}_{\mathrm{K}}(\mathrm{A})\) . Then \(\xi\) is effective if and only if \(u(\xi)\) is.

Lemma 3. --- Let S and T be two nonisomorphic simple A-modules that are finite-dimensional over K. The supports of the \(\mathrm{A}_{(\mathrm{L})}\) -modules \(\mathrm{S}_{(\mathrm{L})}\) and \(\mathrm{T}_{(\mathrm{L})}\) are disjoint.

By Lemma 2, there exists an element \(a \in \mathrm{A}\) such that \(a_{\mathrm{S}} = 0\) and \(a_{\mathrm{T}} = 1\) . The element \(1 \otimes a\) of \(\mathrm{A}_{(\mathrm{L})}\) acts as 0 on \(\mathrm{S}_{(\mathrm{L})}\) and as 1 on \(\mathrm{T}_{(\mathrm{L})}\) . By Lemma 2, the supports of the \(\mathrm{A}_{(\mathrm{L})}\) -modules \(\mathrm{S}_{(\mathrm{L})}\) and \(\mathrm{T}_{(\mathrm{L})}\) are disjoint.

Let us prove Theorem 1. Let S be the set of classes of simple A-modules that are finite-dimensional over K. The family \(([S])_{S \in \mathcal{S}}\) is a basis of the Z-module \(R_{K}(A)\) . Let \(S \in S\) . The \(A_{(L)}\) -module \(S_{(L)}\) is not zero, so its support is not empty. Let \(S'\) be an element of this support. By Lemma 1 of VIII, p.~190, there exists a homomorphism \(f_{S'} : R_L(A_{(L)}) \to Z\) such that \(f_{S'}([E]) = \ell_{S'}(E)\) for every \(A_{(L)}\) -module E that is finite-dimensional over L. We have \(f_{S'}([S_{(L)}]) \neq 0\) by construction and \(f([T_{(L)}]) = 0\) for every \(T \in S - \{S\}\) by Lemma 3. We have therefore proved that the elements of \(R_L(A_{(L)})\) of the form \([S_{(L)}]\) for \(S \in S\) are linearly independent over Z. It follows that the homomorphism u is injective.

Let \(S \in \mathscr{S}\) , and let \(S'\) be an element of the support of \(S_{(L)}\) . For every \(\xi \in R_K(A)\) , the coordinate of \(\xi\) of index {[}S{]} in the basis \(([S])_{S \in \mathscr{S}}\) is \(f_{S'}(u(\xi)) / [S_{(L)} : S']\) . It follows that if \(u(\xi)\) is effective, then so is \(\xi\) .

